\chapter{Discussion}
\label{ch:discussion}

% Written 21 September 2026 from analysis/SPINE.md sections 2, 7b and 8.

Chapter~\ref{ch:results} reported a null on the pre-registered arms, and two
findings that the null does not contain: chained cross-episode credit
extinguished the resource on every seed at \(m=2\), and the exploitative
outcome associated with opponent shaping appeared at \(m=3\) under a partner
with no shaping objective.
This chapter says what those results mean, what they do not mean, and what
they imply for opponent shaping between language-model agents.

\section{What the experiment establishes}
\label{sec:disc-establishes}

The environment was built so that a shaper would have something to buy.
At \(m=2\) no player can keep the pool alive alone, commitment loses to
reciprocity, and a restrained partner is worth 55.1 units to a best
responder (Chapter~\ref{ch:env}).
Before any arm ran, two gates established that this stack can acquire the
relevant behaviour here: restraint against a committed harvester, and
reciprocity against tit-for-tat (Section~\ref{sec:gates}).
In that environment, with the objective isolated from its update schedule by
a trial-batched control at the same learning rate, clip, batch and prompt,
the trial-level objective did not move the co-learner's acquired policy in
the direction the claim requires.

This is a negative result about a specific estimator of a specific
objective, obtained in a setting constructed so that a positive result was
detectable.
That is the sense in which it is informative: the design's own manipulation
checks passed, the controls acquired the target behaviour, and the arms that
carried the objective did not.

\section{The estimator, not the objective, carried the difference}
\label{sec:disc-estimator}

The sharpest comparison in the experiment is internal.
Shaper-matched and the split-credit arm differ in nothing except how
cross-episode credit is estimated: the same seat, learning rate, clip,
update schedule, batch and prompt, and in both cases the shaper is trained
on the return of the whole trial.
Under chained whole-trial GAE the resource was extinguished on every seed
and the co-learner ended mixed or none.
Under credit split into within-episode advantages plus a later-episode
return taken against a baseline from previous trials, the resource survived
in 0.58 of episodes, the co-learner became unconditional on four of five
seeds, and the frozen policy transferred to a fresh learner on five of five.
Neither arm beat the controls; one matched them and the other destroyed the
thing being contested.

The mechanism is measured, not inferred (Section~\ref{sec:res-v1}).
The estimator as first executed bootstrapped the value through the
environment's reset at every collapse, and on the chained arms' own
records its push toward restraint reverses sign, on every seed, when that
one bootstrap is set to zero.
Chaining through a collapse into a fresh pool inverts the signal, and the
inversion feeds itself once collapses are common.
Split credit computes advantages within episodes and never bootstraps
through a reset, which is why it sustained the resource with the same
objective, seat, speed, schedule, batch and prompt.
Three further properties of the chained estimator are real and secondary:
credit reaches a step discounted by \(\lambda\) per live step, so a reward
one episode later enters the opening advantage of a full episode at
\(0.97^{50}=0.22\); whitening per update removes the component of the
cross-episode term that is common to the trial, which is the component a
shared co-learner's response produces; and in the pilot diagnosis the
whole-trial advantages had a standard deviation of 14 to 32 against the
naive learner's 2 to 7.
% [Corrected chained arms: one sentence on what they did, from the rerun.]

The practical implication is narrow and useful.
An opponent-shaping objective for language-model agents is not only a choice
of objective; it is a choice of estimator, and at this scale the estimator
dominates the outcome.
Reporting the objective without reporting how its credit is assigned leaves
the result unidentified.

\section{Where the exploitative outcome came from}
\label{sec:disc-m3}

At \(m=3\) the ShapeLLM-style arm produced the pattern the shaping
literature reports: agent~2 earning 1.81 units per live round against
agent~1's 1.01, with the asymmetric profile on 0.83 of rounds.
Reported alone, in the currency that opponent-shaping papers use, that is a
successful shaping result.
It is not one here, for a reason fixed before the data: the slow arm, whose
agent~2 differs from an ordinary naive learner only in learning rate and
clip, produced the same signature, and decision~C passed on three of three
seeds.

The mechanism is not mysterious.
At \(m=3\) a single player can sustain the pool and commitment wins
(Chapter~\ref{ch:env}), and an agent that updates more slowly is closer to a
committed player.
A slow partner therefore collects the asymmetric profile without any
cross-episode credit at all.
This is the two-timescale reading of the outcome rather than a shaping
reading, and the design separates them because the regimes differ in one
environment parameter and nothing else.

The consequence for how such results should be read is general.
In games where commitment pays, an exploitative asymmetry between a slower
agent and a faster one is the expected outcome of the speed difference.
Evidence that an objective produced it therefore requires a control that
holds the speed and the update schedule fixed --- which is what the
trial-batched arms supply here, and what the published comparison, whose
only shaping control is an observation-space ablation, does not.

\section{Relation to ShapeLLM}
\label{sec:disc-shapellm}

This thesis asks ShapeLLM's question in a different environment class, and
its result should not be read as a refutation of that paper.
Three differences bound the comparison.
The environment is a sequential resource game with a stock and a horizon of
50 rounds, not a repeated \(2\times2\) matrix game at \(T=20\).
The hyperparameters are this stack's, and they differ from the paper's in
GAE \(\lambda\), value coefficient, mini-batch, score scaling, entropy and
the shaper's learning rate and clip (Chapter~\ref{ch:design}).
And the readout is the co-learner's acquired policy under a scripted probe,
not the reward each agent collects while both are still training.

What transfers is the identification argument.
The paper's evidence is average reward per step and joint-action visitation
for a shaper that also learns ten times more slowly and updates five times
less often, in games that include Chicken, where commitment pays.
The present experiment shows that in a game of that family the same evidence
form is produced by a slow partner with no shaping objective.
The claim that the objective is responsible therefore needs the schedule and
timescale controls, independently of whether the claim is true.
That is a statement about what the evidence can identify, not about whether
LLM agents can shape.

\section{Implications for shaping on shared resources}
\label{sec:disc-commons}

One result deserves emphasis beyond the shaping question.
The arms that optimised the shaper's trial return with the estimator as
first executed drove a renewable resource to extinction in every episode of
every seed, earning 0.38 to 0.49 units per round where mutual harvest pays
0.36 and mutual restraint pays 1.00, and the cause was a bootstrap through
the environment's reset that no matrix-game benchmark would have exposed.
An objective that is sound on repeated matrix games can, through its
estimator, produce worse outcomes for both agents than not optimising it
at all when the shared object can be exhausted.

The reading is the first of the two that were possible: the estimator
pointed the wrong way, and Section~\ref{sec:res-v1} shows the sign
reversing on the recorded play once the bootstrap through the reset is
removed.
The alternative --- that the trial return is maximised by harvesting hard
while the partner still restrains --- is what the corrected chained arms
test, and their result is read in Section~\ref{sec:res-mechanism}.
% [From the rerun: whether the corrected arms sustained the resource.]

For deployment the weaker statement is already enough.
An agent trained to maximise its own return across an interaction in which a
co-learner adapts can arrive at a policy that is worse for itself and for
the other party than an agent trained on the immediate episode.
Where the shared object is a resource that can be exhausted, the failure
mode is not a lost bargain but a destroyed commons.

\section{What would change the answer}
\label{sec:disc-next}

The result is bounded by one base model, one adapter rank, one environment
family, 100 epochs and three to five seeds (Chapter~\ref{ch:lim}).
Within those bounds, three specific things would move it.

A lower-variance estimator of the same trial objective is the first.
The split-credit arm shows that the estimator changes the outcome; it does
not show that a better one produces shaping, because it matched the controls
rather than exceeding them.

A co-learner that does not acquire the target behaviour on its own is the
second.
At \(m=2\) the naive baseline became unconditional on every seed, so the
design's headroom for a shaping effect was smaller than the solver's stake
suggested.
An environment, prior or budget in which the control does not already
succeed would restore it --- at the cost of a new pre-registration, since
changing the design point after seeing these results would convert a
registered null into an unregistered positive.

A direct test of whether the chained estimator can learn restraint at all
against a stationary, reciprocal partner is the third, and it is cheap: it
requires no co-learner and therefore cannot produce a shaping result, only
locate the failure.
It is specified in Chapter~\ref{ch:lim}.
